\section{Lecture 12: evaluation 的核心问题}

前面课程已经讲了如何训练语言模型：架构、训练、系统、scaling、推理。Lecture 12 转向一个更基础也更容易被低估的问题：给定一个模型，如何判断它“好不好”？这个问题必须先于数据选择，因为数据会塑造模型行为；如果我们不知道目标行为是什么，就不知道该训练什么数据。

\begin{importantbox}{Evaluation 的核心挑战}
Evaluation 不是机械地“给 prompts、拿 responses、算 accuracy”。它是在把一个抽象 construct，例如“能力”“有用性”“安全性”“推理能力”“真实世界表现”，压缩成 concrete metric。这个压缩过程会反过来塑造研究方向、产品选择和社会判断。
\end{importantbox}

\begin{knowledgebox}{课堂提示：评价先问目的，再选指标}
老师在本讲开头的主线不是“介绍一串 benchmark 名字”，而是提醒我们：there is no one true evaluation。评价要服务具体问题。用户或公司想做购买决策，研究者想测 raw capability，政策制定者想理解风险和收益，模型开发者想拿反馈改模型；这四种目的需要不同 eval，混用会导致错误结论。
\end{knowledgebox}

\begin{knowledgebox}{术语消化：评价设计的三条轴}
\begin{tabular}{L{0.22\textwidth}L{0.34\textwidth}L{0.34\textwidth}}
\toprule
轴 & 问题 & 例子 \\
\midrule
Difficulty & 任务是否仍能区分 frontier models？ & MMLU 被 MMLU-Pro/HLE 替代。 \\
Realism & 是否像真实用户或真实工作？ & Chatbot Arena、GDPVal、MedHELM。 \\
Validity & 指标是否真的测到想测的 construct？ & contamination、dataset quality、judge bias。 \\
\bottomrule
\end{tabular}
\end{knowledgebox}

\subsection{本章小结}

本讲的入口问题不是“哪个榜单最权威”，而是“你想用这个榜单回答什么问题”。后面所有指标都要回到这个入口：它测什么、漏什么、被什么污染、会怎样塑造模型开发。

